\chapter{Uncertainty, Ehrenfest's Theorem, and the Classical Limit}

\section{The Loose End in the Linear Hamiltonian}

The model
\begin{equation}
H=cP
\label{eq:l10-linear-hamiltonian}
\end{equation}
was useful because its wave packets translate without changing shape. It also carries a serious warning. Momentum may be positive or negative, so the same formula assigns positive energy to one plane wave and negative energy to the other:
\begin{equation}
\psi_p(x)=e^{ipx/\hbar},
\qquad
E(p)=cp.
\end{equation}
If arbitrarily many physical particles could occupy negative-energy states, empty space would not be the state of lowest energy. It could produce positive-energy quanta while depositing compensating negative energy into new particles, preserving total energy but lowering the vacuum without bound.

\begin{classroomqa}
\audiencequestion{Why is a negative-energy branch more than an odd choice of zero for the energy?}
\lecturerresponse{A common additive shift of every energy is harmless. An unbounded spectrum is different. If ever lower-energy states can be populated while positive-energy particles carry away the balance, there is no stable ground state. Stability requires a lower bound relative to the vacuum.}
\end{classroomqa}

\begin{figure}[htbp]
\centering
\includegraphics[width=0.88\textwidth]{lecture_10_figure_01.png}
\caption{The linear Hamiltonian and its plane-wave eigenfunctions expose the negative-energy branch when \(p<0\).}
\lectureframe{00:04:35}
\end{figure}

Dirac's historical answer worked for fermions. Fermions obey the exclusion principle: no two identical fermions can occupy the same one-particle state. Imagine filling every negative-energy state. Once all of them are occupied, exclusion prevents another fermion from falling into that sector. A vacancy in the filled set behaves as an antiparticle.

This argument does not rescue negative-energy bosons. Any number of bosons may occupy one state, so a negative-energy bosonic mode could be populated without limit. The old filled-sea picture is best regarded as the historical route to antiparticles; modern quantum field theory organizes particles and antiparticles without treating the physical vacuum as a literal sea of observable negative-energy electrons. The point needed here is narrower: the toy Hamiltonian is not a complete stable particle theory.

As a one-dimensional caricature of a fermionic particle moving at a fixed speed, the model hints at the neutrino problem for which Dirac's ideas were relevant. It cannot serve as a photon theory: photons are bosons, and exclusion cannot fill and block their negative-energy states.

\begin{classroomqa}
\audiencequestion{Why can exclusion stabilize the fermionic construction but not a bosonic one?}
\lecturerresponse{A fermionic state can be filled once and then blocked. A bosonic state has no such occupancy ceiling, so a negative-energy bosonic mode would still let the energy fall without bound.}
\end{classroomqa}

\paragraph{Where the course now stands.}
The quarter has developed much of the foundation of quantum mechanics and several deliberately simple systems, but very few applications. In particular, it has not yet reached even the harmonic oscillator. That omission marks a boundary rather than an ending: relativity and field theory come next, after which quantum mechanics must be taken up again to study the oscillator and the other structures built on it. The sustained interest in the earlier discussion of entanglement also shows why the foundational route was worthwhile.

With that perspective, the loose end is closed. The main business is the wave function, uncertainty, and the relation between quantum evolution and classical mechanics. The aim is not the most general possible evolution equation, but the special position-space equation originally written by Schr\"odinger.

\section{Wave Functions and Complete Commuting Coordinates}

For one particle on a line, the position-space wave function is
\begin{equation}
\psi(x)=\braket{x}{\psi},
\end{equation}
the amplitude for finding the particle near \(x\). The momentum-space wave function is
\begin{equation}
\widetilde{\psi}(p)=\braket{p}{\psi}.
\end{equation}
They describe the same abstract state in different bases and are related by a Fourier transform:
\begin{align}
\psi(x)
&=
\int_{-\infty}^{\infty}
\frac{dp}{\sqrt{2\pi\hbar}}\,
\widetilde{\psi}(p)e^{ipx/\hbar},
\label{eq:l10-fourier-forward}\\
\widetilde{\psi}(p)
&=
\int_{-\infty}^{\infty}
\frac{dx}{\sqrt{2\pi\hbar}}\,
\psi(x)e^{-ipx/\hbar}.
\label{eq:l10-fourier-inverse}
\end{align}

\begin{figure}[htbp]
\centering
\includegraphics[width=0.88\textwidth]{lecture_10_figure_02.png}
\caption{Position and momentum wave functions are projections of one state onto two different continuous bases.}
\lectureframe{00:09:05}
\end{figure}

In two spatial dimensions, \(X\) and \(Y\) commute, so a common eigenbasis may be labeled by both coordinates and the wave function becomes \(\psi(x,y)\). That commutation is not forced by the abstract axioms alone; it is part of the empirically successful quantum description of ordinary space. For \(N\) particles in three dimensions,
\begin{equation}
\psi(\mathbf{x}_1,\ldots,\mathbf{x}_N)
\end{equation}
is a function on configuration space, not a separate three-dimensional wave attached independently to each particle. A macroscopic system can therefore require an enormous number of coordinates.

\begin{classroomqa}
\audiencequestion{Does the commutation of different spatial coordinates follow automatically from quantum mechanics?}
\lecturerresponse{No. The abstract framework permits many algebras of observables. The fact that distinct Cartesian position components commute is an empirical structural fact about the quantum mechanics of particles in ordinary space.}
\end{classroomqa}

One may Fourier transform every coordinate, obtaining a wave function of all momenta, or transform only selected variables. Mixed representations are legitimate because a representation is a basis choice, not a change in the physical state.

\section{What Uncertainty Measures}

The most primitive form of the uncertainty principle follows at once from the Fourier pair. A state of definite momentum is a plane wave with constant modulus and no position localization. A state sharply localized in position requires many momentum components. More generally, narrowing one member of a Fourier pair broadens the other.

\begin{figure}[htbp]
\centering
\includegraphics[width=0.88\textwidth]{lecture_10_figure_03.png}
\caption{A localized position profile requires a broad superposition of plane waves in the reciprocal momentum representation.}
\lectureframe{00:14:45}
\end{figure}

To make that statement quantitative, begin with a general observable \(Q\). Its expectation value in a normalized state is
\begin{equation}
\expect{Q}=\bra{\psi}Q\ket{\psi}.
\end{equation}
The variance is
\begin{equation}
(\Delta Q)^2
=
\expect{\bigl(Q-\expect{Q}\bigr)^2}
=
\expect{Q^2}-\expect{Q}^2.
\label{eq:l10-general-variance}
\end{equation}
The uncertainty \(\Delta Q\) is the positive square root of this variance.

For position one may first shift the origin so that \(\expect{X}=0\). Then
\begin{equation}
(\Delta x)^2
=
\expect{X^2}
=
\int_{-\infty}^{\infty}
dx\,|\psi(x)|^2x^2.
\label{eq:l10-position-variance}
\end{equation}
Positive and negative contributions cancel in \(\expect{X}\), but they add in \(\expect{X^2}\). A broader probability distribution therefore has a larger variance.

\begin{figure}[htbp]
\centering
\includegraphics[width=0.88\textwidth]{lecture_10_figure_04.png}
\caption{With the mean shifted to zero, the position variance is the probability-weighted integral of \(x^2\).}
\lectureframe{00:20:50}
\end{figure}

The distinction between \(\expect{X^2}\) and \(\expect{X}^2\) is essential. A distribution symmetric about zero has \(\expect{X}=0\), but unless it is concentrated at the origin it still has \(\expect{X^2}>0\).

Momentum variance may be defined directly in momentum space:
\begin{equation}
(\Delta p)^2
=
\int_{-\infty}^{\infty}
dp\,|\widetilde{\psi}(p)|^2
\bigl(p-\expect{P}\bigr)^2.
\end{equation}
It is more useful for the proof to work in position space. Shift the mean momentum to zero and use
\begin{equation}
P=-i\hbar\frac{d}{dx},
\qquad
P^2=-\hbar^2\frac{d^2}{dx^2}.
\end{equation}
Then
\begin{equation}
(\Delta p)^2
=
\expect{P^2}
=
-\hbar^2
\int_{-\infty}^{\infty}
dx\,\psi^*(x)\frac{d^2\psi}{dx^2}.
\label{eq:l10-momentum-variance-second}
\end{equation}

\begin{classroomqa}
\audiencequestion{How can a variance be positive when its position-space formula begins with a minus sign?}
\lecturerresponse{The integral multiplying that sign is nonpositive. Integrating by parts transfers one derivative to \(\psi^*\), introduces a second minus sign, and turns the result into the integral of an absolute square.}
\end{classroomqa}

For wave functions that decay sufficiently rapidly at infinity,
\begin{equation}
\int f\,g'\,dx
=
\left[fg\right]_{-\infty}^{\infty}
-
\int f'g\,dx,
\end{equation}
and the boundary term vanishes. Applying this once to Equation~\eqref{eq:l10-momentum-variance-second} gives
\begin{equation}
(\Delta p)^2
=
\hbar^2
\int_{-\infty}^{\infty}
dx\,
\frac{d\psi^*}{dx}
\frac{d\psi}{dx}
=
\hbar^2
\int_{-\infty}^{\infty}
dx\,
\left|\frac{d\psi}{dx}\right|^2
\geq 0.
\label{eq:l10-momentum-variance-positive}
\end{equation}

\begin{figure}[htbp]
\centering
\includegraphics[width=0.88\textwidth]{lecture_10_figure_05.png}
\caption{Integration by parts converts the apparent negative expression for \(\expect{P^2}\) into a manifestly positive norm.}
\lectureframe{00:27:35}
\end{figure}

\begin{classroomqa}
\audiencequestion{Can the means of both position and momentum be shifted to zero at the same time?}
\lecturerresponse{Yes. Translating the argument of the wave function shifts \(\expect{X}\), while multiplying it by a plane-wave phase shifts \(\expect{P}\). These operations alter the means without changing the corresponding variances, so both centered variables may be used in the proof.}
\end{classroomqa}

\section{The Uncertainty Bound}

The mathematical tool needed is the Cauchy--Schwarz inequality. The lecture approaches it through elementary vector geometry: the dot product of two vectors is their lengths multiplied by the cosine of the angle between them, so its magnitude cannot exceed the product of their lengths. In Hilbert-space notation,
\begin{equation}
\braket{A}{A}\braket{B}{B}
\geq
\left|\braket{A}{B}\right|^2.
\label{eq:l10-cauchy-schwarz}
\end{equation}
It holds in real or complex vector spaces of any dimension.

\begin{figure}[htbp]
\centering
\includegraphics[width=0.88\textwidth]{lecture_10_figure_06.png}
\caption{The geometric vector inequality is recast as the Hilbert-space Cauchy--Schwarz bound used in the proof.}
\lectureframe{00:35:35}
\end{figure}

\begin{classroomqa}
\audiencequestion{Should the inequality be strict, or may equality occur?}
\lecturerresponse{It is a greater-than-or-equal relation. Equality occurs when the two Hilbert-space vectors are linearly dependent. That equality condition later identifies the minimum-uncertainty wave packets.}
\end{classroomqa}

Center both variables and first follow the simplified blackboard proof, which takes \(\psi(x)\) to be real. Define
\begin{equation}
\ket{A}=X\ket{\psi},
\qquad
\ket{B}=P\ket{\psi}.
\label{eq:l10-auxiliary-vectors}
\end{equation}
Their squared norms are precisely the variances:
\begin{equation}
\braket{A}{A}=(\Delta x)^2,
\qquad
\braket{B}{B}=(\Delta p)^2.
\end{equation}

\begin{figure}[htbp]
\centering
\includegraphics[width=0.88\textwidth]{lecture_10_figure_07.png}
\caption{The auxiliary vectors \(X\ket{\psi}\) and \(P\ket{\psi}\) turn the two variances into Hilbert-space norms.}
\lectureframe{00:44:40}
\end{figure}

The remaining inner product is
\begin{align}
\braket{A}{B}
&=
\int_{-\infty}^{\infty}
dx\,
x\psi(x)
\left(-i\hbar\frac{d\psi}{dx}\right).
\end{align}
Only its absolute value enters. Since \(\psi\) is real,
\begin{equation}
\psi\frac{d\psi}{dx}
=
\frac{1}{2}\frac{d}{dx}\psi^2.
\end{equation}
Therefore,
\begin{align}
\left|\braket{A}{B}\right|
&=
\frac{\hbar}{2}
\left|
\int_{-\infty}^{\infty}
dx\,
x\frac{d}{dx}\psi^2
\right| \notag\\
&=
\frac{\hbar}{2}
\left|
-\int_{-\infty}^{\infty}
dx\,\psi^2
\right|
=
\frac{\hbar}{2},
\label{eq:l10-inner-product-half-hbar}
\end{align}
where integration by parts removes the derivative from \(\psi^2\), \(d x/dx=1\), the boundary term vanishes, and normalization supplies the last equality.

Insert Equation~\eqref{eq:l10-inner-product-half-hbar} into Cauchy--Schwarz:
\begin{equation}
(\Delta x)^2(\Delta p)^2
\geq
\frac{\hbar^2}{4}.
\end{equation}
Taking the positive square root gives
\begin{equation}
\boxed{
\Delta x\,\Delta p
\geq
\frac{\hbar}{2}
}.
\label{eq:l10-uncertainty-principle}
\end{equation}

\begin{figure}[htbp]
\centering
\includegraphics[width=0.88\textwidth]{lecture_10_figure_08.png}
\caption{The integration-by-parts step and normalization produce the sharp position-momentum uncertainty bound.}
\lectureframe{00:50:55}
\end{figure}

The restriction to real \(\psi\) was a blackboard convenience, not a restriction on the theorem. For a general complex state, take
\begin{equation}
\ket{A}=(X-\expect{X})\ket{\psi},
\qquad
\ket{B}=(P-\expect{P})\ket{\psi}.
\end{equation}
The commutator fixes the imaginary part of their overlap:
\begin{equation}
\operatorname{Im}\braket{A}{B}
=
\frac{1}{2i}\expect{[X,P]}
=
\frac{\hbar}{2}.
\end{equation}
Since \(|\braket{A}{B}|\) is at least as large as the magnitude of its imaginary part, the same bound follows.

\begin{classroomqa}
\audiencequestion{What exactly happened to the boundary term in the repeated integrations by parts?}
\lecturerresponse{For the normalizable packets considered here, the products appearing at \(+\infty\) and \(-\infty\) vanish. On a finite periodic interval, the endpoint terms cancel. The domain of the operators must be chosen so that these statements hold; otherwise the formal manipulation is not justified.}
\end{classroomqa}

\begin{classroomqa}
\audiencequestion{Is the equality sign physically attainable, or is the lower bound merely an estimate?}
\lecturerresponse{It is attainable. Equality in Cauchy--Schwarz requires the two auxiliary vectors to be proportional. Solving that condition produces Gaussian packets, the minimum-uncertainty states developed next.}
\end{classroomqa}

\section{Other Fourier-Conjugate Pairs}

The uncertainty relation between time and energy resembles the one between position and momentum. It requires more care, however, because time is usually an evolution parameter here, not an observable represented by a universal self-adjoint operator. Expand a state in energy eigenvectors:
\begin{equation}
\ket{\psi(t)}
=
\int dE\,
\phi(E)e^{-iEt/\hbar}\ket{E}.
\label{eq:l10-time-energy-fourier}
\end{equation}
Energy amplitudes and temporal variation are Fourier related. A narrow energy distribution changes slowly and cannot make a sharply timed pulse; a broad energy distribution can produce a short arrival-time envelope.

\begin{figure}[htbp]
\centering
\includegraphics[width=0.88\textwidth]{lecture_10_figure_09.png}
\caption{Time dependence is assembled from energy eigenstates with phases \(e^{-iEt/\hbar}\), giving the Fourier basis of the time-energy relation.}
\lectureframe{00:53:40}
\end{figure}

\begin{classroomqa}
\audiencequestion{Does the same proof apply directly to time and energy?}
\lecturerresponse{The reciprocal-width intuition does, and a passing wave packet gives a useful example: a short arrival-time window requires a broad energy distribution. The operator statement is subtler because time does not enter this formulation exactly as \(X\) does.}
\end{classroomqa}

Angle and angular momentum form another conjugate pair for rotational motion. Locally their Fourier relation resembles position and momentum. Globally the periodicity of angle introduces additional care, so the remark is an analogy here rather than a full theorem.

\begin{classroomqa}
\audiencequestion{Are position-momentum, time-energy, and angle-angular momentum isolated examples?}
\lecturerresponse{No. Canonically conjugate variables repeatedly appear as Fourier pairs. The details of domains and periodicity differ, but reciprocal localization is the common structure.}
\end{classroomqa}

\section{The Schr\"odinger Equation in Position Space}

We now return to dynamics. The abstract law is
\begin{equation}
i\hbar\frac{d}{dt}\ket{\psi(t)}
=
H\ket{\psi(t)}.
\label{eq:l10-abstract-schrodinger}
\end{equation}
For a particle of mass \(m\) in a potential \(V(X)\), take
\begin{equation}
H=\frac{P^2}{2m}+V(X).
\label{eq:l10-particle-hamiltonian}
\end{equation}
This imports the classical kinetic-plus-potential form and then tests whether its quantum evolution reproduces classical motion in the proper regime.

\begin{figure}[htbp]
\centering
\includegraphics[width=0.88\textwidth]{lecture_10_figure_10.png}
\caption{The Hamiltonian acts on the abstract state vector in the time-dependent Schr\"odinger equation.}
\lectureframe{00:59:05}
\end{figure}

In the position representation, \(P^2\) becomes a second derivative. A function of \(X\) acts by multiplication:
\begin{equation}
\bigl(V(X)\psi\bigr)(x)=V(x)\psi(x).
\end{equation}
Consequently,
\begin{equation}
i\hbar\frac{\partial\psi(x,t)}{\partial t}
=
\left[
-\frac{\hbar^2}{2m}\frac{\partial^2}{\partial x^2}
+V(x)
\right]\psi(x,t).
\label{eq:l10-position-schrodinger}
\end{equation}

\begin{figure}[htbp]
\centering
\includegraphics[width=0.88\textwidth]{lecture_10_figure_11.png}
\caption{The particle Hamiltonian becomes the kinetic second derivative plus multiplication by \(V(x)\).}
\lectureframe{01:03:05}
\end{figure}

\begin{classroomqa}
\audiencequestion{In what sense is the ordinary function \(V(x)\) an operator?}
\lecturerresponse{Once a position basis is chosen, \(X\) acts by multiplying by \(x\). A sufficiently regular function \(V(X)\) therefore acts by multiplying by the corresponding number \(V(x)\) at each point. Its usefulness is tested by the dynamics that follows.}
\end{classroomqa}

For the expectation-value calculations, solve Equation~\eqref{eq:l10-position-schrodinger} for the two time derivatives:
\begin{align}
\dot{\psi}
&=
\frac{i\hbar}{2m}\psi''
-\frac{i}{\hbar}V\psi,
\label{eq:l10-psi-dot}\\
\dot{\psi}^{\,*}
&=
-\frac{i\hbar}{2m}\psi^{*\,\prime\prime}
+\frac{i}{\hbar}V\psi^*.
\label{eq:l10-psi-star-dot}
\end{align}
The second line is only the complex conjugate of the first, but writing both makes the cancellations visible.

\section{The Motion of the Packet Center}

The center of the position distribution is
\begin{equation}
\expect{X}
=
\int_{-\infty}^{\infty}
dx\,\psi^*(x,t)\,x\,\psi(x,t).
\end{equation}
Differentiate it:
\begin{equation}
\frac{d}{dt}\expect{X}
=
\int dx\,
\left(
\dot{\psi}^{\,*}x\psi
+
\psi^*x\dot{\psi}
\right).
\label{eq:l10-x-expectation-derivative}
\end{equation}
The two terms are complex conjugates. On substituting Equations~\eqref{eq:l10-psi-dot} and \eqref{eq:l10-psi-star-dot}, the potential terms cancel because \(V(x)\) is real. Integrating the kinetic terms by parts leaves
\begin{align}
\frac{d}{dt}\expect{X}
&=
\frac{1}{m}
\int dx\,
\psi^*
\left(-i\hbar\frac{\partial}{\partial x}\right)
\psi \notag\\
&=
\frac{\expect{P}}{m}.
\label{eq:l10-ehrenfest-x}
\end{align}

\begin{figure}[htbp]
\centering
\includegraphics[width=0.88\textwidth]{lecture_10_figure_12.png}
\caption{The derivative of the packet center reduces to the expectation value of momentum divided by the mass.}
\lectureframe{01:15:40}
\end{figure}

This relation is exact. No narrow-packet approximation entered the algebra.

\begin{classroomqa}
\audiencequestion{Where was the assumption that the wave packet is a single narrow bump used in this derivation?}
\lecturerresponse{Nowhere. The identity \(d\expect{X}/dt=\expect{P}/m\) holds for every state in the domain of this Hamiltonian. Localization matters only when we interpret \(\expect{X}\) as the trajectory of a classical particle.}
\end{classroomqa}

The other half of mechanics concerns momentum:
\begin{equation}
\expect{P}
=
\int dx\,
\psi^*
\left(-i\hbar\frac{\partial}{\partial x}\right)
\psi.
\end{equation}
Differentiate both factors that depend on time:
\begin{equation}
\frac{d}{dt}\expect{P}
=
-i\hbar
\int dx\,
\left(
\dot{\psi}^{\,*}\psi'
+
\psi^*\frac{\partial\dot{\psi}}{\partial x}
\right).
\label{eq:l10-p-expectation-start}
\end{equation}
After integrating one term by parts, substitute the two Schr\"odinger equations. The kinetic contributions cancel. The potential terms combine into
\begin{equation}
\int dx\,
V(x)
\left(
\psi\frac{\partial\psi^*}{\partial x}
+
\psi^*\frac{\partial\psi}{\partial x}
\right)
=
\int dx\,
V(x)\frac{\partial}{\partial x}|\psi|^2.
\end{equation}
One final integration by parts moves the derivative onto the potential:
\begin{align}
\frac{d}{dt}\expect{P}
&=
-\int dx\,|\psi(x,t)|^2\frac{dV}{dx} \notag\\
&=
-\expect{V'(X)}
\\
&=
\expect{F(X)},\\
F(x)&=-V'(x).
\label{eq:l10-ehrenfest-p}
\end{align}

\begin{figure}[htbp]
\centering
\includegraphics[width=0.88\textwidth]{lecture_10_figure_13.png}
\caption{The potential terms combine into a total derivative, yielding \(d\expect{P}/dt=-\expect{V'(X)}\).}
\lectureframe{01:30:45}
\end{figure}

Equations~\eqref{eq:l10-ehrenfest-x} and \eqref{eq:l10-ehrenfest-p} are Ehrenfest's theorem for this Hamiltonian. They are exact quantum identities, but exact expectation-value equations are not yet the same as a classical orbit.

\section{Where the Classical Approximation Enters}

Classical mechanics for a point at the packet center would require
\begin{equation}
\frac{d}{dt}\expect{P}
\approx
F\!\left(\expect{X}\right).
\label{eq:l10-classical-center-force}
\end{equation}
Quantum mechanics gives instead
\begin{equation}
\frac{d}{dt}\expect{P}
=
\expect{F(X)}.
\end{equation}
In general,
\begin{equation}
\expect{F(X)}
\neq
F\!\left(\expect{X}\right).
\label{eq:l10-force-nonlinearity}
\end{equation}

\begin{classroomqa}
\audiencequestion{If both Ehrenfest equations are exact, why do they not already prove that every quantum state follows Newton's law?}
\lecturerresponse{Newton's law for the packet center needs the force evaluated at the center. The exact quantum equation contains the average of the force over the full probability distribution. Those coincide for a narrow packet in a slowly varying force field, but not for an arbitrary state or nonlinear force.}
\end{classroomqa}

The lecture's counterexample makes the distinction unmistakable. Let \(F(x)=x^2\) and choose a normalized state with two equal narrow bumps centered near \(+a\) and \(-a\). Symmetry gives
\begin{equation}
\expect{X}=0,
\qquad
F\!\left(\expect{X}\right)=0,
\end{equation}
while
\begin{equation}
\expect{F(X)}
=
\expect{X^2}
\approx
a^2>0.
\end{equation}

\begin{figure}[htbp]
\centering
\includegraphics[width=0.88\textwidth]{lecture_10_figure_14.png}
\caption{For a symmetric two-bump packet and \(F(x)=x^2\), the force at the mean vanishes while the mean force does not.}
\lectureframe{01:33:35}
\end{figure}

For a single narrow packet centered at \(\bar{x}=\expect{X}\), expand the force:
\begin{equation}
F(X)
=
F(\bar{x})
+
F'(\bar{x})(X-\bar{x})
+
\frac{1}{2}F''(\bar{x})(X-\bar{x})^2
+\cdots.
\end{equation}
Taking the expectation value removes the linear term:
\begin{equation}
\expect{F(X)}
=
F(\bar{x})
+
\frac{1}{2}F''(\bar{x})(\Delta x)^2
+\cdots.
\end{equation}
The classical approximation is reliable when the packet remains narrow compared with the length scale on which the force changes appreciably.

Sharp features can split and scatter a packet. A broad, smooth potential tends to guide a sufficiently localized packet without destroying its single-bump form; a pointlike obstacle, atomic core, or structure on the scale of the wavelength can produce reflected and diffracted components.

\begin{figure}[htbp]
\centering
\includegraphics[width=0.88\textwidth]{lecture_10_figure_15.png}
\caption{Potential structure on a scale shorter than the packet width or relevant wavelength can break a coherent packet into scattered components.}
\lectureframe{01:37:25}
\end{figure}

\begin{classroomqa}
\audiencequestion{What makes a potential bad for a classical packet?}
\lecturerresponse{The issue is scale, not moral character. Features much sharper than the packet or comparable to its wavelength resolve the wave nature and generate scattering. If the potential varies slowly across the packet, the force is nearly constant over its support and the center follows the classical equation.}
\end{classroomqa}

Mass also matters. For a nearly minimum-uncertainty packet,
\begin{equation}
\Delta x\,\Delta v
\sim
\frac{\hbar}{2m}.
\end{equation}
A large mass permits both a small spatial width and a modest velocity spread. Heavy objects can therefore remain sharply localized while traversing macroscopic potentials. Light particles are more readily spread or diffracted by microscopic structure. This is a scaling argument, not a claim that mass alone decides whether motion is classical.

\begin{classroomqa}
\audiencequestion{Is an electron ever large enough, in this sense, to behave approximately classically?}
\lecturerresponse{Yes, if the external field varies smoothly on the electron packet's scale. Between broad capacitor plates, for example, a coherent packet can follow the classical acceleration for a useful time. Near the sharp Coulomb structure of a nucleus, the same electron displays pronounced scattering.}
\end{classroomqa}

Rutherford's 1911 scattering experiment provides the complementary lesson. Alpha particles encountering the compact charge distribution of a gold nucleus were scattered through large angles, revealing that the positive charge was concentrated in a very small region. Rutherford described a beam of particles; quantum mechanics describes the incident amplitude and its scattered components. Modern electron and photon scattering experiments use the same scale sensitivity.

\begin{classroomqa}
\audiencequestion{Is this wave-versus-scale behavior special to matter particles?}
\lecturerresponse{No. Electromagnetic waves follow geometric optics when the refractive index changes slowly over a wavelength. A diffraction grating with spacing comparable to the wavelength spreads the wave into distinct directions. Electrons, alpha particles, photons, and other quantum excitations share this relation between wavelength and structure.}
\end{classroomqa}

The classical limit is therefore not a separate law pasted onto quantum mechanics. The exact expectation-value equations already contain the classical relations. Classical trajectories emerge when the state remains localized and the Hamiltonian varies smoothly enough that averaging the force is effectively the same as evaluating it at the packet center.
